% conclusion_summary.tex — un-headed synthesis under \section{Conclusion}

Every factor analysis of a rating scale rests on two structures: the
covariance of responses and the organization of the language that
elicited them. The first has a century of method, the second no unified toolkit until
now; \texttt{semanticfa} brings the full discipline of exploratory factor
analysis---explicit encoding choices, consensus retention, matched-null
diagnostics, interpretation, and refinement---to that second structure,
respondent-free and exactly reproducible from fixed embeddings. On a
well-understood inventory it yielded a \valConsensusKWord-factor recovery of the Big
Five; congruence with the \valNHuman-respondent structure strong for
Openness, Conscientiousness, and Neuroticism, informatively weaker where
English merges warmth and sociability; anti-topic vectors, not
reverse-scored meanings, from sign-flipped reverse-keyed embeddings; and
pre-data refinement that flagged redundant items, sharpened a short form,
and vetted candidates. Agreement between the two
structures---semantic--behavioral coherence---belongs in the nomological
network of any construct measured in language: convergence shows words and
respondents organizing the construct alike; divergence, where they do not
and why. Neither replaces the other; semantic factor analysis estimates,
before any respondent is recruited, how much of any scale's structure is
already written in its items.
